%% ============================================================
%%  Lecture 2 -- Thursday, 28 May 2026
%%  Subfile of main.tex: compiles standalone OR as part of the book.
%%  Built from sources/2026-05-28/ (board-transcript.md [Gate-1 verified], outline.md).
%% ============================================================
\documentclass[main.tex]{subfiles}

\begin{document}

\beginlecture{2}{Thursday, 28 May 2026}%
  {Counting the Infinite, and the Geometry of Distance: Countability and Metric Spaces}

\epigraph{No one shall expel us from the paradise that Cantor has created.}{David Hilbert}

\begin{lecturecontents}{Contents of this lecture}
  \item \hyperlink{2.1}{2.1\quad Counting the infinite: countable and uncountable sets}
  \item \hyperlink{2.2}{2.2\quad Metric spaces}
  \item \hyperlink{2.3}{2.3\quad Open sets and open balls}
\end{lecturecontents}

\bigskip
\noindent\textbf{Overview.} Lecture~1 built $\mathbb R$ as the complete ordered field and drew from
completeness the Archimedean property (\xrefL{1.7.1}) and the density of $\mathbb Q$ (\xrefL{1.7.2});
the same property gives infima as readily as suprema (\xrefL{1.5.3}). Two new questions open here.
First, \emph{how large} is $\mathbb R$? Cantor's notion of countability lets us compare infinite
sets, and his diagonal argument shows $\mathbb R$ is strictly larger than $\mathbb Q$---uncountable,
with $\mathbb Q$ a vanishingly ``thin'' subset. Second, what is the right language for
\emph{distance}? Metric spaces distil the Euclidean norm into three axioms, and the open sets they
generate are the gateway to topology, where the rest of the course lives.
The companion text is Rudin, Chapter~2, especially \R{2.1}--\R{2.19}.

\clearpage
% ============================================================
\sub{2.1}{Counting the infinite: countable and uncountable sets}

\begin{Obj}{Definition}{2.1.1}{Injective, surjective, bijective}
Let $f:X\to Y$ be a function (cf.\ \R{2.1}--\R{2.3}).
\begin{itemize}
  \item $f$ is \emph{injective} if $f(x)=f(x')\Rightarrow x=x'$ (distinct inputs give distinct outputs);
  \item $f$ is \emph{surjective} if every $y\in Y$ equals $f(x)$ for some $x\in X$;
  \item $f$ is \emph{bijective} if it is both. A bijection has a two-sided inverse $f^{-1}:Y\to X$.
\end{itemize}
\end{Obj}

\begin{Obj}{Definition}{2.1.2}{Equinumerous sets; finite, countable, uncountable}
Sets $X$ and $Y$ are \emph{equinumerous}, written $X\sim Y$, if there is a bijection $f:X\to Y$
(\R{2.3}). With
$J_0=\emptyset$ and $J_n=\{1,\dots,n\}$, a set $X$ is
\begin{itemize}
  \item \emph{finite} if $X\sim J_n$ for some $n$, and \emph{infinite} otherwise;
  \item \emph{countable} if $X\sim\mathbb N$;
  \item \emph{uncountable} if it is infinite and $X\not\sim\mathbb N$.
\end{itemize}
This packages Rudin's finite/countable terminology \R{2.4}--\R{2.7}.
\end{Obj}

\begin{Obj}{Proposition}{2.1.3}{Equinumerosity is an equivalence relation}
The relation $\sim$ is reflexive, symmetric, and transitive. This is the basic relation behind
\R{2.3}.
\end{Obj}
\begin{Prf}{2.1.3}{Direct Proof that Equinumerosity is an Equivalence Relation.}{[Direct]\quad bijections invert and compose, with $(g\circ f)^{-1}=f^{-1}\circ g^{-1}$.}
\item For reflexivity, the identity $\mathrm{id}_X:X\to X$, $x\mapsto x$, is a bijection, so $X\sim X$.
\item For symmetry, if $X\sim Y$ via a bijection $f$, then $f^{-1}:Y\to X$ is a bijection, so $Y\sim X$.
\item For transitivity, if $X\sim Y$ via $f$ and $Y\sim Z$ via $g$, then $g\circ f:X\to Z$ is a
  bijection (its inverse is $f^{-1}\circ g^{-1}$), so $X\sim Z$.
\end{Prf}

\begin{ObjL}{Example}{2.1.4}{$\mathbb Z$ is countable}
The identity shows $\mathbb N$ is countable. For $\mathbb Z$, interleave the signs:
\[ f:\mathbb N\to\mathbb Z,\qquad f(n)=\begin{cases} n/2 & n\text{ even},\\[2pt] -(n+1)/2 & n\text{ odd},\end{cases} \]
a bijection that lists $\mathbb Z$ as $0,-1,1,-2,2,-3,3,\dots$ (cf.\ \R{2.12}--\R{2.13}).
\end{ObjL}

\begin{Fig}{2.1.5}{$\mathbb N$ counts $\mathbb Z$}{The bijection of \xref{2.1.4} walks outward from $0$ in alternating steps, reaching every integer exactly once.}
\begin{tikzpicture}[x=1cm,y=1cm]
  \node[left,cuBlueDk] at (-0.55,1.3){$\mathbb N$};
  \node[left,cuBlueDk] at (-0.55,0){$\mathbb Z$};
  \foreach \n/\fz in {0/0,1/{-1},2/1,3/{-2},4/2,5/{-3},6/3}{
    \node[figTerm] (a\n) at (\n,1.3){\scriptsize$\n$};
    \node (b\n) at (\n,0){\scriptsize$\fz$};
    \draw[->,cuBlueDk!70] (a\n)--(b\n);
  }
\end{tikzpicture}
\end{Fig}

\begin{Obj}{Definition}{2.1.6}{Sequences and the listing criterion}
A \emph{sequence} in a set $X$ is a function $f:\mathbb N\to X$, written $\{x_n\}_{n\in\mathbb N}$ with
$x_n=f(n)$. Thus $X$ is countable precisely when some sequence lists every element of $X$ exactly once.
Rudin's sequence definition is \R{2.7}.
\end{Obj}

\begin{ObjL}{Example}{2.1.7}{$\mathbb Q$ is countable}
Identifying a fraction with its pair of integers exhibits $\mathbb Q$ inside
$\mathbb Z\times\mathbb Z$, which a diagonal sweep lists as a single sequence.
\end{ObjL}
\begin{Prf}{2.1.7}{Proof by Diagonal Enumeration that $\mathbb Q$ is Countable.}{[Constr]\quad cf.\ \R{2.12}--\R{2.13}.}
\item Arrange the pairs $(p,q)$ with $q\neq0$ on a grid and traverse them along successive finite
  diagonals (each diagonal $|p|+|q|=\text{const}$ holds only finitely many pairs).
\item This runs through every pair; deleting pairs with a common factor and any repeats leaves a
  sequence that lists each rational exactly once.
\item Hence $\mathbb Q\sim\mathbb N$ (\xref{2.1.6}).
\end{Prf}

\begin{Fig}{2.1.8}{Listing the rationals}{Sweeping the lattice of pairs $(p,q)$ along finite diagonals threads them into one sequence; pairs with a common factor (e.g.\ $\tfrac22,\tfrac24$) are skipped (\xref{2.1.7}).}
\begin{tikzpicture}[x=1.0cm,y=0.9cm]
  \foreach \q in {1,2,3,4}{\foreach \p in {1,2,3,4}{\filldraw[cuBlueDk] (\q,\p) circle (1pt);}}
  \node[cuBlueDk,below=1pt] at (2.5,0.6){\scriptsize denominator $q$};
  \node[cuBlueDk,rotate=90] at (0.45,2.5){\scriptsize numerator $p$};
  \draw[->,figTerm,thick] (1,1)--(1,2)--(2,1)--(3,1)--(2,2)--(1,3)--(1,4)--(2,3)--(3,2)--(4,1);
  \foreach \q/\p in {2/2,2/4,4/2,3/3,4/4}{\node[figBad] at (\q,\p){\scriptsize$\times$};}
\end{tikzpicture}
\end{Fig}

\begin{Obj}{Proposition}{2.1.9}{An infinite subset of a countable set is countable}
If $X$ is countable and $A\subseteq X$ is infinite, then $A$ is countable (\R{2.8}).
\end{Obj}
\begin{Prf}{2.1.9}{Proof by Recursive Selection of Indices that the Subset is Countable.}{[A10]\quad cf.\ \R{2.8}; the recursion is an induction, justified at each step by $A$ being infinite.}
\item Write $X=\{x_n\}_{n\in\mathbb N}$, every element listed once (\xref{2.1.6}).
\item Choose indices recursively: let $n_0$ be the least index with $x_{n_0}\in A$; having chosen
  $n_0<\cdots<n_k$, let $n_{k+1}$ be the least index $>n_k$ with $x_{n_{k+1}}\in A$. Such an index
  exists because $A$, being infinite, cannot lie inside the finite set $\{x_0,\dots,x_{n_k}\}$, so
  some $x_m\in A$ has $m>n_k$.
\item The map $k\mapsto x_{n_k}$ is strictly increasing, hence injective, and reaches every element
  of $A$ (each $a=x_m\in A$ is selected no later than stage $m$); so it is a bijection
  $\mathbb N\to A$, and $A\sim\mathbb N$.
\end{Prf}

\begin{Obj}{Theorem}{2.1.10}{Cantor: the binary sequences are uncountable}
The set $2^{\mathbb N}=\{\,f:\mathbb N\to\{0,1\}\,\}$ of infinite binary sequences---equivalently, the
subsets of $\mathbb N$---is uncountable.
\end{Obj}
\begin{Prf}{2.1.10}{Proof by Cantor's Diagonal Argument that the Binary Sequences are Uncountable.}{[A8$+$A9]\quad diagonalisation, by contradiction; cf.\ \R{2.14}.}
\item Suppose $2^{\mathbb N}$ were countable, listed as $a_1,a_2,a_3,\dots$, each
  $a_k=(a_k^1,a_k^2,a_k^3,\dots)$ a binary sequence.
\item Define $\bar b\in 2^{\mathbb N}$ by flipping the diagonal: $\bar b_n=1-a_n^{\,n}$.
\item For each $n$, $\bar b$ differs from $a_n$ in its $n$th entry ($\bar b_n\neq a_n^{\,n}$), so
  $\bar b\neq a_n$ for every $n$.
\item Thus $\bar b$ appears nowhere in the list---contradicting that the list was all of
  $2^{\mathbb N}$. Hence $2^{\mathbb N}$ is uncountable.
\end{Prf}

\begin{Fig}{2.1.11}{Cantor's diagonal}{Flipping the boxed diagonal entries yields a sequence $\bar b$ differing from each listed $a_k$ in place $k$, so $\bar b$ is absent from the list (\xref{2.1.10}).}
\begin{tikzpicture}[x=0.66cm,y=0.6cm]
  \foreach \row/\digits/\yy in {1/{0,1,1,0,1,0}/5, 2/{1,1,1,1,0,1}/4, 3/{0,0,0,1,0,1}/3,
                                4/{0,1,1,0,1,0}/2, 5/{1,0,0,0,0,1}/1}{
    \node[cuBlueDk] at (0,\yy){\scriptsize$a_{\row}$};
    \foreach \d [count=\c] in \digits {\node at (\c,\yy){\scriptsize$\d$};}
  }
  \foreach \c/\yy in {1/5,2/4,3/3,4/2,5/1}{\draw[figBad,thick] (\c,\yy) circle (3.4pt);}
  \draw[figBad] (0.7,5.3)--(5.3,0.7);
  \node[figBad] at (0,-0.3){\scriptsize$\bar b$};
  \foreach \d [count=\c] in {1,0,1,1,1}{\node[figBad] at (\c,-0.3){\scriptsize$\d$};}
  \node at (6,-0.3){\scriptsize$\cdots$};
\end{tikzpicture}
\end{Fig}

\begin{Obj}{Theorem}{2.1.12}{The interval $[0,1]$ is uncountable}
The interval $[0,1]\subseteq\mathbb R$ is uncountable.
\end{Obj}
\begin{Prf}{2.1.12}{Proof by Reduction to Binary Sequences that $[0,1]$ is Uncountable.}{[Direct]\quad uses \xref{2.1.10}; barring an all-$1$ tail makes binary expansions unique ($0.0111\ldots_2=0.1000\ldots_2$), and the barred eventually-all-$1$ sequences are countable.}
\item Since $[0,1]\supseteq[0,1)$, it suffices to show $[0,1)$ is uncountable. Send each $x\in[0,1)$
  to the binary expansion $(b_n)\in 2^{\mathbb N}$ that does \emph{not} end in all $1$s; on $[0,1)$
  this representative exists and is unique, so the map is injective, with image exactly the
  sequences not ending in all $1$s.
\item If $[0,1)$ were countable, that image would be countable, so $2^{\mathbb N}$---the image
  together with the countably many all-$1$-tail sequences---would be countable too (\R{2.12}),
  contradicting \xref{2.1.10}.
\item Hence $[0,1)$, and therefore $[0,1]$, is uncountable.
\end{Prf}

\begin{Obj}{Corollary}{2.1.13}{$\mathbb R$ is uncountable}
Were $\mathbb R$ countable, its infinite subset $[0,1]$ would be countable too (\xref{2.1.9}); but
it is not (\xref{2.1.12}). So $\mathbb R$ is uncountable.
Compare Rudin's diagonal proof for binary sequences \R{2.14} and later perfect-set result \R{2.43}.
\end{Obj}

\begin{ObjL}{Remark}{2.1.14}{$\mathbb Q$ is thin in $\mathbb R$}
Thus $\mathbb Q$ is countable yet dense (\xrefL{1.7.2}), while $\mathbb R$ is uncountable: the
rationals are a vanishingly small part of the line. Made precise by measure, $\mathbb Q$---indeed
any countable set---has Lebesgue length zero, so a real drawn from any distribution with a density
is irrational with probability $1$. Compare the countability results in \R{2.12}--\R{2.13} and
the diagonal uncountability result \R{2.14}; the measure language is later perspective.
\end{ObjL}

% ============================================================
\sub{2.2}{Metric spaces}

\noindent The Euclidean distance on $\mathbb R^n$, $d(x,y)=\|x-y\|=\big(\sum_{i=1}^n|x_i-y_i|^2\big)^{1/2}$,
built from the inner product $\langle x,y\rangle=\sum_i x_iy_i$, is the model. A metric space keeps
only the features of $d$ that analysis actually uses.

\begin{Obj}{Definition}{2.2.1}{Metric space}
A \emph{metric space} (\R{2.15}) is a set $X$ with a \emph{metric} (or distance)
$d:X\times X\to\mathbb R^{\ge0}$ satisfying, for all $x,y,z\in X$:
\begin{enumerate}
  \item $d(x,y)\ge0$, with $d(x,y)=0\iff x=y$;
  \item $d(x,y)=d(y,x)$ \quad(symmetry);
  \item $d(x,z)\le d(x,y)+d(y,z)$ \quad(the triangle inequality).
\end{enumerate}
\end{Obj}

\begin{Fig}{2.2.2}{The triangle inequality}{The direct distance $d(x,z)$ never exceeds the detour $d(x,y)+d(y,z)$ through $y$---axiom~(3) of \xref{2.2.1}.}
\begin{tikzpicture}[x=1cm,y=1cm]
  \coordinate (x) at (0,0); \coordinate (y) at (4,0); \coordinate (z) at (1.4,2.1);
  \draw[figTerm,thick] (x)--(z) node[midway,above left]{$d(x,z)$};
  \draw[cuBlueDk] (x)--(y) node[midway,below]{$d(x,y)$};
  \draw[cuBlueDk] (y)--(z) node[midway,above right]{$d(y,z)$};
  \foreach \p/\nm/\pos in {x/x/{below left}, y/y/{below right}, z/z/above}{
    \filldraw (\p) circle (1.3pt) node[\pos]{$\nm$};}
\end{tikzpicture}
\end{Fig}

\begin{ObjL}{Example}{2.2.3}{The standard metrics}
$(\mathbb R,|\cdot|)$ with $d(x,y)=|x-y|$ is a metric space, as are $(\mathbb Z,|\cdot|)$,
$(\mathbb N,|\cdot|)$, and $(\mathbb Q,|\cdot|)$. More generally any subset $Y$ of a metric space
$(X,d)$ is itself a metric space under the restricted distance $d|_{Y\times Y}$. The Euclidean space
$(\mathbb R^n,d_2)$, $d_2(x,y)=\big(\sum_{i=1}^n|x_i-y_i|^2\big)^{1/2}$, is the prototype.
These are Rudin's basic metric examples (cf.\ \R{2.16}).
\end{ObjL}

\begin{ObjL}{Example}{2.2.4}{The $p$-metrics and the discrete metric}
On $\mathbb R^n$, for $p\ge1$,\note{The board writes $p>1$; in fact $d_p$ is a metric for every $p\ge1$, with $d_1$ and $d_\infty$ the boundary cases.}
\[ d_p(x,y)=\Big(\sum_{i=1}^n|x_i-y_i|^p\Big)^{1/p},\qquad d_\infty(x,y)=\max_i|x_i-y_i|. \]
On any set $X$, the \emph{discrete metric}---$d(x,y)=0$ if $x=y$ and $1$ otherwise---is a metric.
The triangle inequality for $d_p$ is Minkowski's inequality, taken here as a standard fact
(cf.\ \R{2.16}).
\end{ObjL}

\begin{ObjL}{Example}{2.2.5}{Metrics on a function space}
Let $C=C[0,1]$ be the continuous functions $f:[0,1]\to\mathbb R$. Both
\[ d_\infty(f,g)=\sup_{[0,1]}|f-g| \qquad\text{and}\qquad d_1(f,g)=\int_0^1|f(x)-g(x)|\,dx \]
are metrics on $C$. Here $d_\infty$ is finite because a continuous function on $[0,1]$ is bounded,
and $d_1(f,g)=0$ forces $f=g$ since $|f-g|$ is continuous. As metric examples, compare \R{2.16}.
\end{ObjL}

% ============================================================
\sub{2.3}{Open sets and open balls}

\noindent Basic \emph{topology} studies a space through its \emph{open subsets}---concepts that can
be phrased using open sets alone. Everything starts with the open ball.

\begin{Obj}{Definition}{2.3.1}{Open ball}
In a metric space $(X,d)$, the \emph{open ball} of radius $r>0$ about $p\in X$ is
\[ B_r(p)=\{\,x\in X:d(x,p)<r\,\}. \]
Rudin calls this a neighbourhood of $p$ in \R{2.18}.
\end{Obj}

\begin{Fig}{2.3.2}{An open ball takes the shape of its metric}{Same centre, same radius: an interval in $(\mathbb R,|\cdot|)$, a round disk under $d_2$, a diamond under $d_1$, a square under $d_\infty$---each the set $\{x:d(x,p)<r\}$ of \xref{2.3.1}.}
\begin{tikzpicture}[x=1cm,y=1cm]
  \begin{scope}[xshift=0cm]
    \draw[cuBlueDk] (-1.05,0)--(1.05,0);
    \draw[figTerm,line width=1.2pt] (-0.8,0)--(0.8,0);
    \draw[figTerm,fill=white] (-0.8,0) circle (1.7pt);
    \draw[figTerm,fill=white] (0.8,0) circle (1.7pt);
    \filldraw (0,0) circle (1.1pt) node[above=2pt]{\scriptsize$p$};
    \node[cuBlueDk] at (0,-1.5){\scriptsize$(\mathbb R,|\cdot|)$};
  \end{scope}
  \begin{scope}[xshift=3.3cm]
    \draw[figTerm,fill=cuBlueLt!35,dashed] (0,0) circle (0.85);
    \draw[cuBlueDk] (0,0)--(0.85,0) node[midway,below=1pt]{\scriptsize$r$};
    \filldraw (0,0) circle (1.1pt) node[above=2pt]{\scriptsize$p$};
    \node[cuBlueDk] at (0,-1.5){\scriptsize$(\mathbb R^2,d_2)$};
  \end{scope}
  \begin{scope}[xshift=6.6cm]
    \draw[figTerm,fill=cuBlueLt!35,dashed] (0,0.88)--(0.88,0)--(0,-0.88)--(-0.88,0)--cycle;
    \filldraw (0,0) circle (1.1pt) node[above right=-1pt]{\scriptsize$p$};
    \node[cuBlueDk] at (0,-1.5){\scriptsize$(\mathbb R^2,d_1)$};
  \end{scope}
  \begin{scope}[xshift=9.9cm]
    \draw[figTerm,fill=cuBlueLt!35,dashed] (-0.82,-0.82) rectangle (0.82,0.82);
    \filldraw (0,0) circle (1.1pt) node[above right=-1pt]{\scriptsize$p$};
    \node[cuBlueDk] at (0,-1.5){\scriptsize$(\mathbb R^2,d_\infty)$};
  \end{scope}
\end{tikzpicture}
\end{Fig}

\begin{Obj}{Definition}{2.3.3}{Interior point; open set}
Let $(X,d)$ be a metric space and $E\subseteq X$. A point $p\in E$ is an \emph{interior point} of $E$
if $B_r(p)\subseteq E$ for some $r>0$, and $E$ is \emph{open in $X$} if every point of $E$ is
interior (\R{2.18}). Conversely, $p$ fails to be interior exactly when $B_r(p)\not\subseteq E$ for
\emph{every} $r>0$; so $E$ is \emph{not} open as soon as one of its points is not interior.
\end{Obj}

\begin{ObjL}{Remark}{2.3.4}{Openness is relative}
Openness is meaningful only \emph{relative to the ambient space}: one says ``$E$ is open in $X$,''
never merely ``$E$ is open.'' The same set can be open in one metric space and not in another.
This anticipates Rudin's relative-openness theorem \R{2.30}.
\end{ObjL}

\begin{Obj}{Proposition}{2.3.5}{Every open ball is open}
In any metric space $(X,d)$, the ball $B_r(p)$ is open in $X$, for every $p\in X$ and $r>0$.
\end{Obj}
\begin{Prf}{2.3.5}{Proof by the Triangle Inequality that Every Open Ball is Open.}{[Direct]\quad the template for openness proofs; uses \xref{2.2.1}; cf.\ \R{2.19}.}
\item Fix $q\in B_r(p)$; then $d(p,q)<r$, so $r_q:=r-d(p,q)>0$.
\item Take any $x\in B_{r_q}(q)$. By the triangle inequality,
  \[ d(x,p)\le d(x,q)+d(q,p)<r_q+d(q,p)=r. \]
\item Hence $x\in B_r(p)$, so $B_{r_q}(q)\subseteq B_r(p)$. Every $q\in B_r(p)$ is therefore interior,
  and $B_r(p)$ is open.
\end{Prf}

\begin{Fig}{2.3.6}{Why a ball is open}{For $q\in B_r(p)$, the radius $r_q=r-d(p,q)$ keeps $B_{r_q}(q)$ inside $B_r(p)$---the triangle inequality at work (proof of \xref{2.3.5}).}
\begin{tikzpicture}[x=1cm,y=1cm]
  \draw[figTerm,thick] (0,0) circle (2);
  \filldraw (0,0) circle (1.2pt) node[below left]{$p$};
  \coordinate (q) at (1.1,0.4);
  \draw[cuBlueDk,dashed] (0,0)--(q) node[midway,above]{\scriptsize$d(p,q)$};
  \draw[figLim,thick] (q) circle (0.8);
  \filldraw (q) circle (1.2pt) node[above right=-1pt]{$q$};
  \draw[cuBlueDk] (q)--(1.9,0.4) node[midway,below=1pt]{\scriptsize$r_q$};
  \node[figTerm] at (1.35,-1.6){$B_r(p)$};
  \node[figLim] at (2.05,1.35){$B_{r_q}(q)$};
\end{tikzpicture}
\end{Fig}

\begin{ObjL}{Remark}{2.3.7}{A recipe for proving openness}
To show $E$ is open in $X$: (1)~take an arbitrary $p\in E$; (2)~produce a radius $r_p>0$ (guided by
the geometry); (3)~verify $B_{r_p}(p)\subseteq E$. The proof of \xref{2.3.5} is the model, and
this is the working form of \R{2.18}--\R{2.19}.
\end{ObjL}

% per-lecture Index of Objects -- standalone compile only; the book uses the master index.
\ifSubfilesClassLoaded{\clearpage\objindex}{}

\end{document}
